% =====================================================================
%  tr06a.tex —— 原文 pp.76--86（tr-06）
%  §13 第五讲（Banach 代数与 C*-代数的 K 理论）：pp.76--81
%  §14 第六讲（C*-代数扩张与 K-同调）：pp.82--86
% =====================================================================
% ===== 原文 pp.76--86 =====
% [原文 p.76]
\section{第五讲（Banach 代数与 C*-代数的 K 理论）}

回忆一下，\emph{Banach 代数}是完备的赋范代数。

$C^*$-代数是指这样的对合 Banach 代数，它满足
\[
  \|xx^*\| = \|x\|^2.
\]

\para{例子：}
\begin{enumerate}[label=\alph*.]
\item Hilbert 空间 $H$ 上有界算子之集 $B(H)$。
\item Hilbert 空间 $H$ 上紧算子之集 $K(H)$（它是 $B(H)$ 的一个理想）。
\item Calkin 代数
  \[
    Q(H) = B(H)/K(H).
  \]
\item 若 $X$ 是紧 Hausdorff 空间，则 $C(X)$（$X$ 上连续复函数之集）赋予上确界范数
  \[
    \|f\| = \sup |f|
  \]
  与对合
  \[
    g^* = \bar g
  \]
  （其中 $\bar g$ 表示复共轭），是一个 $C^*$-代数。
\end{enumerate}

% [Tutor Note: p.76 $C^*$-代数的范数条件原文排作 $\|xx^*\|=\|x\|^2$，这里没有笔误：由对合等距（$\|x^*\|=\|x\|$）可得 $\|x^*x\|=\|xx^*\|=\|x\|^2$，所以它与教材常见的 $\|x^*x\|=\|x\|^2$ 等价。核对方法：把 $x$ 换成 $x^*$，并用对合的对合性 $(x^*)^*=x$ 即可。]
% [Tutor Note: p.76 例 c 的 Calkin 代数 $Q(H)=B(H)/K(H)$ 与 p.78 提到的 Quillen $+$ 构造没有关系（作者在 p.78 特意提醒区分两个「$+$」）。$Q(H)$ 是第六讲的主线：p.80 的边界映射 $\partial$ 在 $J=K(H)$、$A=B(H)$ 时就是广义指标，p.82 Berg 定理用 $\sigma(\pi(N))$ 定义本质谱，p.83 $BDF(X)$ 又等于单同态 $C(X)\to Q(H)$ 之集。]
% [Tutor Note: p.76 例 d 的显示式原文写作 $g^*=\bar g$，而紧邻上一行用的是 $f$（$\|f\|=\sup|f|$），原文记号不统一，译文照抄页图。按上下文它表示 $C(X)$ 的对合 $f^*(x)=\overline{f(x)}$，即逐点取复共轭。]

\para{旁注 1：}若 $X$ 只是局部紧且 Hausdorff，则考虑 $C_0(X)$（在无穷远处消失的复连续函数之集）。它也是无单位的 $C^*$-代数。此外
\[
  C_0(X) = \operatorname{Ker}(C(X^+)\to\C),
\]
其中
\[
  f\mapsto f(\infty),
\]
% [原文 p.77]
这里 $X^+$ 是 $X$ 的一点紧化。

每个非单位的 $C^*$-代数都可以按这种方式「1-紧化」：
\[
  J \to J^+ \to \C,
\]
其中
\[
  J^+ = \{j+\lambda 1 : \lambda\in\C\}.
\]

\para{旁注 2：}由 Gelfand 定理我们知道，紧 Hausdorff 空间 $X$ 的拓扑由 $C(X)$（它是交换 $C^*$-代数）所刻画，因此每个交换 $C^*$-代数都对应一个紧 Hausdorff 空间（它的谱）。于是\emph{拓扑 K 理论就是交换 $C^*$-代数的 K 理论}，从而对一般 $C^*$-代数的研究对应于\emph{非交换拓扑}。这正是 A. Connes 著名的\emph{非交换几何}的出发点。

% [Tutor Note: p.77 旁注 2 的「对应」是反变等价：$X\mapsto C(X)$ 与 $A\mapsto\operatorname{Spec}(A)$（$A$ 的特征标空间 / 极大理想空间）互为逆函子（Gelfand--Naimark 定理）。可核对的两点：$X$ 紧 Hausdorff 时 $C(X)$ 的谱恰是 $X$；单位元对应常函数 $1$。正因为如此，p.79 才敢把「对 $(X,A)$ 的正合列」直接翻译成 $0\to C_0(U)\to C(X)\to C(A)\to 0$。]

现在设 $A$ 是有单位的 Banach 代数，$P_n(A)$ 记 $M_n(A)$（元素取自 $A$ 的 $n\times n$ 矩阵之集）中投影（$P^2=P$）之集。用 $P(A)$ 记极限
\[
  P(A) = \lim_{n\to\infty} P_n(A).
\]

\begin{defn}{1}
\[
  K_0 A = \mathrm{Gr}[\pi_0(P(A))],
\]
其中 $\mathrm{Gr}$ 是半群 $\pi_0(P(A))$ 的 Grothendieck 群。
\end{defn}

\begin{defn}{2}
\[
  K_1 A = \pi_0(GL A).
\]
\end{defn}

两个等价的定义如下：

\begin{defn}{1$'$}
我们用生成元与关系把 $K_0(A)$ 定义为如下阿贝尔群：
\begin{itemize}
\item 对每个投影 $p\in M_n(A)$（$\forall n\in\N$）有一个生成元 $[p]$。
\item 关系如下：
  \begin{enumerate}[label=\roman*.]
  \item $[0]=0$
  \item 若 $p,q$ 可通过投影同伦，则 $[p]=[q]$。
% [原文 p.78]
  \item $[p\oplus q]=[p]+[q]$。
  \end{enumerate}
\end{itemize}
\end{defn}

\begin{defn}{2$'$}
我们用生成元与关系把 $K_1(A)$ 定义为如下阿贝尔群：
\begin{itemize}
\item 对每个酉元 $u\in M_n(A)$ 有一个生成元 $[u]$（$\forall n\in\N$，其中酉元指 $u^*u=uu^*=1$）。
\item 关系如下：
  \begin{enumerate}[label=\roman*.]
  \item $[1]=0$
  \item 若 $u,v$ 可通过酉元同伦，则 $[u]=[v]$。
  \item $[u\oplus v]=[u]+[v]$。
  \end{enumerate}
\end{itemize}
\end{defn}

% [Tutor Note: p.77 定义 1 与定义 1$'$ 的等价性核对点：关系 iii（$[p\oplus q]=[p]+[q]$）正是 Grothendieck 群的加法，而关系 ii 把 $\pi_0(P(A))$ 的连通分支（同伦类）当作元素；$\mathrm{Gr}$ 只对交换半群取群化，故 $\pi_0(P(A))$ 必须是交换的，其加法由分块对角 $\operatorname{diag}(p,q)\in M_{2n}(A)$ 经 $\varinjlim$ 给出。这也是为什么先要对 $n$ 取极限 $P(A)=\varinjlim P_n(A)$ 再取 $\pi_0$。]

\begin{lem}{1}
\[
  K_0 A = K_0^{\mathrm{alg}} A,
\]
即我们上面给出的第 0 个 $K$ 群的定义与代数第 0 个 $K$ 群的定义相同。此外在 $P(A)$ 中，同伦与共轭生成相同的等价关系。
\end{lem}

但
\[
  K_1 A \neq K_1^{\mathrm{alg}} A,
\]
因为显然
\[
  \pi_0(GL A) \neq \pi_1(BGL A^+).
\]

% [Tutor Note: p.78 引理 1 的「同伦与共轭生成相同等价关系」补证思路：$C^*$-代数中若两个投影 $p,q\in M_n(A)$ 由投影路径 $p_t$ 相连，可构造 $A$ 中可逆元 $u$ 使 $q=upu^{-1}$（沿路径取 $u_t$ 并注意 $p_t=u_tpu_t^{-1}$）；反之共轭当然同伦。对应到模论：$p\mapsto pA^n$ 把投影的同构类换成有限生成投射模的同构类（Serre--Swan），于是 $K_0^{\mathrm{alg}}$ 与 $K_0$ 一致。]
% [Tutor Note: p.78 的 $K_1A\neq K_1^{\mathrm{alg}}A$ 是本讲最易误读的一处：代数 K 理论中 $K_1^{\mathrm{alg}}A=GL(A)^{\mathrm{ab}}=\pi_1(BGL A^+)$（Quillen $+$ 构造，见 §20），而这里的解析 $K_1A=\pi_0(GL A)$ 是拓扑群的连通分支群。两者一般不同：前者是 $GL(A)=\varinjlim GL_n(A)$ 的阿贝尔化，后者是其连通分支数。$\pi_1(BGL A^+)$ 与 $\pi_0(GL A)$ 之间没有必然的等号，故原文的 $\neq$ 成立。]

若 $J$ 是非单位的 $C^*$-代数，用 $J^+$ 记 $J$ 添上单位后所得的代数；于是我们如下定义 $J$ 的 $K$ 群：
\[
  K_i J = \operatorname{Ker}(K_i(J^+)\to K_i\C), \quad i = 0,1.
\]

我们希望读者能区分一点紧化的「$+$」与 Quillen $+$ 构造的「$+$」。

\emph{基本关系：}
\begin{enumerate}[label=\arabic*.]
\item
  \[
    K_i(C(X)) = K^{-i}X.
  \]
\item 若 $K(H)$ 记 Hilbert 空间 $H$ 上紧算子之集，则
  \[
    K_0(K(H)) = \Z
  \]
% [原文 p.79]
  而
  \[
    K_1(K(H)) = 0.
  \]
  回忆
  \[
    K(H) = \lim_{n\to\infty} M_n\C,
  \]
  于是
  \[
    K_*(K(H)) = \lim_{n\to\infty} K_*(M_n\C) = \lim_{n\to\infty} K_*\C = K_*\C.
  \]
\item 若 $B(H)$ 记 Hilbert 空间 $H$ 上有界算子之集，则
  \[
    K_i(B(H)) = 0, \quad i = 0,1.
  \]
  此外还有
  \[
    K_0(\C) = \Z
  \]
  而
  \[
    K_1(\C) = 0.
  \]
\end{enumerate}

接下来我们问：\emph{对应于对 $(X,A)$ 的正合列，$C^*$-代数 K 理论中的版本应该是什么？}

设 $U = X/A$。则存在正合列
\[
  0 \to C_0(U) \to C(X) \to C(A) \to 0.
\]
有如下结果

\begin{thm}{1}
对任一 $C^*$-代数的短正合列
\[
  0 \to J \to A \to A/J \to 0,
\]
都相伴一个 K 理论群的正合列
\[
  \begin{tikzcd}
    K_1 J \arrow[r] & K_1 A \arrow[r] & K_1(A/J) \arrow[d,"\partial"] \\
    K_0(A/J) \arrow[u] & K_0 A \arrow[l] & K_0 J \arrow[l]
  \end{tikzcd}
\]
% [原文 p.80]
其中竖箭头
\[
  K_0(A/J) \to K_1 J
\]
与
\[
  \partial : K_1(A/J) \to K_0 J
\]
使上图交换。
\end{thm}

% [Tutor Note: p.79 定理 1 的六项列方向务必按图核对：上排从左到右是 $K_1J\to K_1A\to K_1(A/J)$，下排从右到左是 $K_0J\to K_0A\to K_0(A/J)$（即图中箭头指向左），左侧竖箭头向上 $K_0(A/J)\to K_1J$，右侧竖箭头向下 $\partial:K_1(A/J)\to K_0J$。串起来是一个循环：$K_1J\to K_1A\to K_1(A/J)\xrightarrow{\partial}K_0J\to K_0A\to K_0(A/J)\to K_1J$。若把下排误读成从左到右，$\partial$ 的出发/到达群就会颠倒。]

\begin{proof}
略，见参考文献。
\end{proof}

设 $u\in M_n(A/J)$ 可逆。则 $\partial[u]$ 应当阻碍把 $u$ 提升为 $M_nA$ 中的可逆元。

我们知道总可以提升
\[
  \begin{pmatrix} u & 0 \\ 0 & u^{-1} \end{pmatrix},
\]
（更确切地说，是它的类）为可逆矩阵 $R\in M_{2n}A$。令 $P\in M_{2n}A$，
\[
  P = \begin{bmatrix} 1 & 0 \\ 0 & 0 \end{bmatrix} \in M_{2n}(J^+)
\]
并令 $Q = R^{-1}PR \in M_{2n}(J^+)$。则
\[
  [P]-[Q] \in K_0 J.
\]

% [Tutor Note: p.80 的 $\partial[u]=[P]-[Q]$ 是「指标 = 投影之差」的标准构造，可逐步核对：$u\in M_n(A/J)$ 可逆，故 $\operatorname{diag}(u,u^{-1})\in M_{2n}(A/J)$ 的类落在 $\pi_0(GL_{2n}(A/J))$ 的稳定像里，可提升为 $R\in M_{2n}A$ 中的可逆元；取 $P=\operatorname{diag}(1,0)$ 与 $Q=R^{-1}PR$，二者都是 $M_{2n}(J^+)$ 中的投影，其差落在 $\operatorname{Ker}(K_0(J^+)\to K_0\C)=K_0J$ 中。当 $J=K(H)$、$A=B(H)$ 时 $[P]-[Q]$ 就是 Fredholm 指标，故原文说 $\partial$ 是广义\emph{指标}。]

一个有趣的特殊情形是：取 $J = K(H)$、$A = B(H)$，于是 $A/J = Q(H)$。此时边界映射 $\partial$ 是一个广义\emph{指标}。

我们接下来回到（$C^*$-代数的）$K$-同调。

\begin{defn}{3}
设 $A = C(X)$，其中 $X$ 是紧的可度量化拓扑空间。Hilbert 空间称为一个\emph{$X$-模}，若存在表示 $\rho : A \to B(H)$。
\end{defn}

\begin{defn}{4}
两个 $X$-模 $H, H'$ 之间的\emph{本质等价}是一个酉算子 $u : H \to H'$，使得
\[
  u\rho(a)u^* - \rho'(a) \in K
\]
% [原文 p.81]
是紧算子。
\end{defn}

\begin{defn}{5}
$X$-模称为\emph{大的}（big），若映射
\[
  \rho : A \to B(H) \to Q(H)
\]
是单射。
\end{defn}

于是有如下重要的

\begin{thm}{2}
（Voiculescu）大模是\emph{吸收的}，即若 $H, H'$ 是两个 $X$-模且 $H$ 是大的，则 $H \oplus H'$ 与 $H$ 本质等价。
\end{thm}

\begin{proof}
见参考文献。
\end{proof}

\begin{cor}{1}
所有大模都彼此本质等价。
\end{cor}

设 $X$ 是某个拓扑空间（紧且可度量化）。选取一个大 $X$-模 $H_X$。

\begin{defn}{6}
$X$ 的\emph{Paschke 对偶} $D(X)$ 是 $C^*$-代数
\[
  D(X) = \{T \in B(H_X) : [T,\rho(a)] \text{ 紧}\}, \quad \forall a \in A = C(X).
\]
\end{defn}

下一步我们定义：

\begin{defn}{7}
\[
  \wt K_i X = (\wt K^i A =) K_{1-i}(D(X)).
\]
\end{defn}

% [Tutor Note: p.81 定义 7 左边用下标 $i$、括号内用上标 $i$（$\wt K_iX=(\wt K^iA=)K_{1-i}(D(X))$），照页图抄写不要「统一」。指标 $1-i$ 来自 Bott 周期性（周期 2）；与 p.84 的关系 $\wt K^{-1}X=K_1(C(X))$ 对照可自检：取 $i=-1$ 得 $K_{2}(D(X))=K_0(D(X))$，即「$K$-同调的一次悬拓与对偶代数的 $K_0$ 相当」。]
% [Tutor Note: p.81 定义 6 中 $[T,\rho(a)]$ 后的 compact 原文是英文单词（数学排版未加空格），译文写作「紧」。注意 $D(X)$ 不是单位代数：它是 $B(H_X)$ 中与 $\rho(C(X))$ 交换模紧算子的算子全体，因此 p.81 末尾的短正合列要用无单位 $C^*$-代数的 $K$ 群（即 p.78 的 $K_iJ=\operatorname{Ker}(K_i(J^+)\to K_i\C)$）。]

在 $K$-同调中存在 $C^*$-代数对 $(X,Y)$ 的正合列：

\begin{defn}{8}
\[
  D(X,Y) = \{T \in D(X) : T\rho(a) \in K(H), \ \forall a \in C_0(X/Y)\}.
\]
\end{defn}

因为 $D(X,Y) \triangleleft D(X)$，得到 $C^*$-代数的短正合列：
\[
  0 \to D(X,Y) \to D(X) \to D(X)/D(X,Y) \to 0.
\]

% [原文 p.82]
\section{第六讲（C*-代数扩张与 K-同调）}

设 $H$ 是 Hilbert 空间，$N \in B(H)$ 是有界算子。

\begin{defn}{1}
$N$ 称为\emph{正规的}，若
\[
  NN^* = N^*N.
\]
\end{defn}

若 $N$ 正规，则 $N$ 生成 $B(H)$ 的一个交换 $C^*$-子代数，事实上可以证明它同构于 $C(\sigma(N))$，其中
\[
  \sigma(N) = \{\lambda \in \C : \nexists (N-\lambda I)^{-1}\}
\]
是 $N$ 的\emph{谱}。

% [Tutor Note: p.82 谱的定义用的是 $\nexists$（「不存在」）：$\sigma(N)=\{\lambda\in\C:\nexists(N-\lambda I)^{-1}\}$ 应读作「$\lambda\in\C$ 使得 $N-\lambda I$ 不可逆」，即 $\lambda$ 是谱点（不是「$N-\lambda I$ 的逆不存在某个值」）。对照紧接的下一段：本质谱是 $\sigma_e(N)=\sigma(\pi(N))$，其中 $\pi:B(H)\to Q(H)$ 是商映射，把 $B(H)$ 换成 Calkin 代数 $Q(H)$ 就得到本质谱，这是 BDF 理论的关键一步。]

\para{任务：}\emph{尝试把正规算子按模紧算子的酉等价分类。}

Berg 证明了这些算子完全由它们的本质谱 $\sigma_e(N)$ 分类，即
\[
  \sigma(\pi(N)), \qquad \pi : B(H) \to Q(H) = B(H)/K(H),
\]
这是 Voiculescu 定理的一个推论。

\begin{defn}{2}
算子 $N \in B(H)$ 称为\emph{本质正规的}，若 $NN^*-N^*N \in K(H)$ 是紧的。
\end{defn}

接下来，Brown、Douglas、Fillmore 问自己如下问题：\emph{本质谱是否是本质正规算子的本质等价的唯一不变量？}

\emph{答案是否定的。}

作为反例，考虑单边移位 $U$：我们知道 $\sigma_e(U) = S^1$ 且 $\mathrm{Ind} U = -1$，但 $\mathrm{Ind} N = 0$，$\forall N$ 正规。

% [Tutor Note: p.82 的单边移位反例是全讲的动机，务必分清两个不变量：本质谱 $\sigma_e$ 与指标 $\mathrm{Ind}$。$U$（单边移位）的 $\sigma_e(U)=S^1$ 与正规算子（如 $S^1$ 上的乘法算子）相同，但 $\mathrm{Ind}U=-1\neq0=\mathrm{Ind}N$（正规算子 Fredholm 指标为 0），故本质谱不足以分类。这正是 p.83 要加入 $\mathrm{Ind}(N-\lambda I)$、p.84 要用 $BDF(X)$ 的原因。注意原文 $S^1$ 的上标是 $1$，不是字母 $l$。]

% [原文 p.83]
对任何不在 $N$ 的本质谱中的 $\lambda$，指标 $\mathrm{Ind}(N-\lambda I) \in \Z$ 是本质等价的不变量。于是我们可以问：若
\begin{enumerate}[label=\arabic*.]
\item $\sigma_e(N_1) = \sigma_e(N_2) = X$,
\item $\mathrm{Ind}(N_1-\lambda_i I) = \mathrm{Ind}(N_2-\lambda_i I)$，对 $\lambda_i \in \C - X$，
\end{enumerate}
那么算子 $N_1$ 与 $N_2$ 是否本质等价？

这个问题的答案是肯定的。

\begin{defn}{3}
我们把集合 $BDF(X)$ 定义为由本质谱为 $X$ 的本质正规算子的本质等价类所组成的集合。
\end{defn}

于是，$BDF(X)$ 等于 $C^*$-代数扩张
\[
  \begin{tikzcd}
    0 \arrow[r] & K(H) \arrow[r] \arrow[d,equal] & A \arrow[r] \arrow[d,"i"] & C(X) \arrow[r] \arrow[d,"i"] & 0 \\
    0 \arrow[r] & K(H) \arrow[r] & B(H) \arrow[r] & Q(H) \arrow[r] & 0
  \end{tikzcd}
\]
之集，其中 $A$ 是由 $N$ 与 $K(H)$ 生成的代数，$i$ 表示包含 $A \xra{i} B(H)$、$C(X) \xra{i} Q(H)$。

\emph{集合 $BDF(X)$ 也等于单同态 $C(X) \to Q(H)$ 之集。}

此外我们知道如下事实：
\begin{itemize}
\item $BDF(X)$ 在直和下是一个半群。
\item 由 Voiculescu 定理，任何正规算子定义一个单位元。
\end{itemize}

\para{例子：}
% [原文 p.84]
\begin{enumerate}[label=\arabic*.]
\item $BDF(I) = 0$，$I = [0,1] \subset \C$。
\item $BDF(S^1) = \Z$。
\end{enumerate}

% [Tutor Note: p.83 的交换图核对点：两行都是正合列，左竖箭头是恒等（图中画作 $\|$，译文用 tikzcd 的 equal 箭头），中、右两个竖箭头都记为 $i$，分别是从 $A$ 到 $B(H)$、从 $C(X)$ 到 $Q(H)$ 的包含映射。顶行即「半分裂扩张」$0\to K(H)\to A\to C(X)\to 0$；$A$ 由 $N$ 与 $K(H)$ 生成，这正是「本质正规算子 $\leftrightarrow$ 扩张」的对应。三种描述（本质等价类 / 扩张 / 单同态 $C(X)\to Q(H)$）在此被说明为同一个集合。]
% [Tutor Note: p.84 例子 2「$BDF(S^1)=\Z$」可自检：$S^1$ 上的单边移位（Toeplitz 算子）给出 $\mathrm{Ind}=-1$ 的类，而 $BDF(S^1)$ 中的指标映射 $\mathrm{Ind}:BDF(S^1)\to\Z$ 是同构，故 $BDF(S^1)\cong\Z$；对照例子 1：$I=[0,1]$ 可缩，所有本质正规算子都本质等价于正规算子，故 $BDF(I)=0$。]

\begin{thm}{1}
（Averson）给定任一同态 $\alpha : C(X) \to Q(H)$，可以找到另一个同态 $\phi : C(X) \to B(H \oplus H)$，使得若
\[
  \phi = \begin{pmatrix} \phi_{11} & \phi_{12} \\ \phi_{21} & \phi_{22} \end{pmatrix},
\]
则
\[
  \pi \circ \phi_{11} = \alpha.
\]

注意 $\phi_{12}, \phi_{21}$ 必须是紧的（即以紧算子为值），而 $\phi_{22}$ 必须是另一个 $*$-同态，记为 $\beta$，因此 $\alpha \oplus \beta$ 可提升到 $B(H)$。
\end{thm}

% [Tutor Note: p.84 定理 1 的署名原文排作「Averson」，应为 \emph{Arveson}（W. B. Arveson，1969：「Notes on extensions of $C^*$-algebras」中的提升定理）。判断依据：这条「任何同态 $C(X)\to Q(H)$ 都可由某个同态 $C(X)\to B(H\oplus H)$ 的 $(1,1)$ 角模紧算子实现」正是 Arveson 提升定理；我已用 500 dpi 裁图复核页图确为「Averson」$a$--$v$--$e$，故正文照原样译出并在此指出。]

\begin{proof}
基本想法如下：
\[
  \phi(xx^*) = \phi(x)\phi(x^*),
\]
\[
  \begin{pmatrix} \phi_{11}(xx^*) & \phi_{12}(xx^*) \\ \phi_{21}(xx^*) & \phi_{22}(xx^*) \end{pmatrix}
  =
  \begin{pmatrix} \phi_{11}(x) & \phi_{12}(x) \\ \phi_{21}(x) & \phi_{22}(x) \end{pmatrix}
  \begin{pmatrix} \phi_{11}(x^*) & \phi_{12}(x^*) \\ \phi_{21}(x^*) & \phi_{22}(x^*) \end{pmatrix},
\]
于是
\[
  \phi_{11}(xx^*) = \phi_{11}(x)\phi_{11}(x^*) + \phi_{12}(x)\phi_{12}(x^*).
\]
% [Tutor Note: p.84 定理 1 证明末式原文印作 $\phi_{11}(xx^*)=\phi_{11}(x)\phi_{11}(x^*)+\phi_{12}(x)\phi_{12}(x^*)$；但按上面的分块矩阵式与矩阵乘法，乘积 $\phi(x)\phi(x^*)$ 的 $(1,1)$ 元应为 $\phi_{11}(x)\phi_{11}(x^*)+\phi_{12}(x)\phi_{21}(x^*)$，即第二项的两个因子都应是 $\phi_{12},\phi_{21}$ 交叉项（$1600$ dpi 裁图复核，页图确实印作 $\phi_{12}(x)\phi_{12}(x^*)$）。正文照原文保留。]
\end{proof}

注意
\[
  BDF(X) \to \Hom(K^{-1}X, \Z)
\]
与
\[
  0 \to K(H) \to A \to C(X) \to 0.
\]
相应的长正合列包含一个映射
\[
  \begin{tikzcd}
    K_1(C(X)) \arrow[r] \arrow[d,equal] & K_0(K(H)) \arrow[d,equal] \\
    K^{-1}X \arrow[r] & \Z
  \end{tikzcd}
\]

% [原文 p.85]
于是 Brown、Douglas、Fillmore 的基本结果是

\begin{thm}{2}
\[
  BDF(X) = K_1 X.
\]
\end{thm}

% [Tutor Note: p.84 末尾图的读法：上排 $K_1(C(X))\to K_0(K(H))$ 是 p.79 定理 1 的长正合列中 $\partial$ 这一段，下排 $K^{-1}X\to\Z$ 用两条竖线（equal 箭头）把上排两个群认作 $K^{-1}X=K_1(C(X))$（p.78「基本关系」1）与 $\Z=K_0(K(H))$（p.78 基本关系 2）。此图把「边界映射」具体化为一个群同态 $K^{-1}X\to\Z$，即 p.84 中 $BDF(X)\to\Hom(K^{-1}X,\Z)$ 的来源。注意此处的分隔符原文用逗号 $\Hom(K^{-1}X,\Z)$，而 p.85、p.86 用分号 $\Hom(K^{-1}X;\Z)$，译文分别照抄。]

\begin{proof}
我们给出证明的梗概：
\[
  \alpha : C(X) \to Q(H)
\]
\[
  \phi : C(X) \to B(H \oplus H),
\]
其中
\[
  \phi = \begin{pmatrix} \phi_{11} & \phi_{12} \\ \phi_{21} & \phi_{22} \end{pmatrix}.
\]
令
\[
  p : H \oplus H \to H
\]
为到第一个因子的投影。则
\[
  [\phi,p] = \begin{pmatrix} 0 & -\phi_{12} \\ \phi_{21} & 0 \end{pmatrix} \in K(H)
\]
恰好是一个 $K$-循环。

一般而言，$X$ 是 $\R^3$ 中的某个多面体，且 $BDF(X) \to \Hom(K^{-1}X; \Z)$。

下面我们继续做一点图追踪：
\[
  \begin{tikzcd}
    K_1 A \arrow[r] \arrow[d] & K_1 X \arrow[r] \arrow[d] & K_1(X,A) \arrow[r] \arrow[d] & K_0 A \arrow[r] \arrow[d] & \null \\
    \Hom(K^{-1}A;\Z) \arrow[r] & \Hom(K^{-1}X;\Z) \arrow[r] & \cdots & \null & \null
  \end{tikzcd}
\]
\end{proof}

% [Tutor Note: p.85 证明里的 $[\phi,p]$ 是分块矩阵的换位子，直接算：$\phi p=\begin{pmatrix}\phi_{11}&0\\ \phi_{21}&0\end{pmatrix}$，$p\phi=\begin{pmatrix}\phi_{11}&\phi_{12}\\ 0&0\end{pmatrix}$，相减得 $\begin{pmatrix}0&-\phi_{12}\\ \phi_{21}&0\end{pmatrix}$；由定理 1 知 $\phi_{12},\phi_{21}$ 紧，故 $[\phi,p]\in K(H)$，这就是一个 Fredholm 模 / $K$-循环。这条换位子给出 $K_1X$ 的类，正是定理 2 的同构 $BDF(X)=K_1X$ 的具体实现。]
% [Tutor Note: p.85 末尾图追踪的两行来源：上排是 $0\to K(H)\to A\to C(X)\to 0$ 的 K 理论长正合列的一段 $K_1A\to K_1X\to K_1(X,A)\to K_0A$；下排是把每个群换成 $\Hom(\wt K^{-1}(-);\Z)$（由 p.84 的 $BDF(X)\to\Hom(K^{-1}X,\Z)$ 与定理 2 得到），竖箭头由定理 2 的同构给出。图中 $\cdots$ 只画到第三列，第四列 $K_0A$ 下的箭头没有落点，原文如此。]

% [原文 p.86]
最后我们得到：
\[
  0 \to \Ext(K^0X;\Z) \to K_1 X \to \Hom(K^{-1}X;\Z) \to 0
\]
以及
\[
  K^{-1}X \to \Z(=K_0(K(H))) \to K_0 A \to K^0 X \to 0(=K_1(K(H))).
\]
我们还用到了关系
\[
  K_1(C(X)) = K^{-1}X
\]
以及
\[
  K_1 X = BDF(X) = \Ext(C(X)) = \Hom(K^{-1}X;\Z).
\]

% [Tutor Note: p.86 第一个正合列是万有系数定理（UCT）在 $K$-同调中的形态：$0\to\Ext(K^0X;\Z)\to K_1X\to\Hom(K^{-1}X;\Z)\to 0$，因 $\Hom$ 项自由故该列分裂。核对方法：与 p.86 第二个正合列对照，后者是 $K^{-1}X\to\Z\to K_0A\to K^0X\to0$，其中 $\Z=K_0(K(H))$、$0=K_1(K(H))$ 用 p.78--p.79 的基本关系 2；两列拼起来就是 p.79 定理 1 的六项列在 $(J,A)=(K(H),A)$ 处的展开。另注意 p.86 第一、二式里 $K$ 的上标是 $0$ 与 $-1$（不是 $K_0X$、$K_1X$ 的下标），页图上 $0$、$-1$ 明显抬高。]
% [Tutor Note: p.86 最后一式把 $K_1X$、$BDF(X)$、$\Ext(C(X))$、$\Hom(K^{-1}X;\Z)$ 并列，末一项按第一式只是中间商／子式的位置，等式应理解为 p.85 已声明的「$X$ 为多面体」情形下的宽松写法，不宜当作对任意紧可度量化 $X$ 都成立的恒等式。术语：$K$-cycle = $K$-循环；semi-group under direct sum = 直和下（构成）半群；Paschke dual = Paschke 对偶（新术语，见回执）。]
